% Осень26!A32; нумерация внутри столбца без пустых строк.
% Источник: лекция 6, §3 (строка 124)
\ticket{29}{Счетная аддитивность интеграла (по множеству)}
\label{ticket:29}

\subsection*{Интеграл как функция множества: счётная аддитивность}

Вопрос: как интеграл зависит от множества интегрирования? Ожидание: при разбиении \(E=E_1\sqcup E_2\) интеграл распадается в сумму. Оказывается, верна и \textbf{счётная} версия.

\begin{lemma}

\emph{ступенчатые функции.} Пусть \(A=\bigsqcup_{k=1}^{\infty}A_k\) --- дизъюнктное разбиение измеримого множества \(A\) на измеримые части (\(k\) пробегает натуральные; не запрещено и конечное множество индексов), и пусть \(g\ge0\) --- ступенчатая (измеримая) функция на \(X\). Тогда \[\int_A g\,d\mu=\sum_{k=1}^{\infty}\int_{A_k} g\,d\mu.\]

\end{lemma}

\begin{proof}

Пусть \(X=\bigsqcup_{i}e_i\) --- допустимое разбиение (конечное!), \(g=\sum_i\lambda_i\chi_{e_i}\), \(\lambda_i\ge0\). По определению интеграла ступенчатой (с учётом сегодняшнего определения интеграла по множеству: \(\int_A g=\int_X g\chi_A\,d\mu\)): \[\int_A g\,d\mu=\sum_i\lambda_i\,\mu(e_i\cap A).\] Каждое \(e_i\cap A\) распадается в дизъюнктное объединение \(e_i\cap A=\bigsqcup_k(e_i\cap A_k)\), и по счётной аддитивности \textbf{меры} \[\mu(e_i\cap A)=\sum_{k=1}^{\infty}\mu(e_i\cap A_k).\] Итого \[\int_A g\,d\mu=\sum_i\sum_k\lambda_i\,\mu(e_i\cap A_k).\] Все слагаемые неотрицательны, поэтому двойную сумму можно переставлять и перегруппировывать как угодно (теорема о перестановке слагаемых положительного ряда / о суммируемых семействах из 2 семестра): \[\sum_i\sum_k\lambda_i\,\mu(e_i\cap A_k)=\sum_k\sum_i\lambda_i\,\mu(e_i\cap A_k)=\sum_{k=1}^{\infty}\int_{A_k}g\,d\mu. \]

\end{proof}

\begin{theorem}

\emph{счётная аддитивность интеграла по множеству, неотрицательный случай.} Пусть \(A=\bigsqcup_{k=1}^{\infty}A_k\) --- дизъюнктное разбиение на измеримые множества и \(f\colon X\to[0,+\infty]\) --- неотрицательная измеримая функция (достаточно измеримости на \(A\)). Тогда \[\int_A f\,d\mu=\sum_{k=1}^{\infty}\int_{A_k} f\,d\mu.\]

\end{theorem}

\begin{proof}

Равенство доказывается «в два приёма»: сначала \(\le\), потом \(\ge\).

\textbf{(≤).} Возьмём любую ступенчатую \(g\) с \(0\le g\le f\) на \(A\). По лемме \[\int_A g\,d\mu=\sum_k\int_{A_k} g\,d\mu\ \le\ \sum_k\int_{A_k} f\,d\mu,\] ибо сужение \(g\) на \(A_k\) остаётся ступенчатым (ступеньки «подрезались» множеством \(A_k\)) и по-прежнему не превосходит \(f\) на \(A_k\). Итак, число \(\sum_k\int_{A_k}f\,d\mu\) --- верхняя грань для всех \(\int_A g\) по ступенчатым \(g\le f\); переходя к супремуму («стандартный приём перехода к супремуму»), \[\int_A f\,d\mu=\sup_g\int_A g\,d\mu\ \le\ \sum_{k=1}^{\infty}\int_{A_k}f\,d\mu.\]

\textbf{(≥), шаг 1: разбиение на два множества} \(A=A_1\sqcup A_2\). Возьмём ступенчатые \(g_1,g_2\) с \[0\le g_1\le f \text{ на } A_1,\qquad 0\le g_2\le f \text{ на } A_2\] и \(g_i=0\) вне \(A_i\). Приведём их к общему допустимому разбиению: \(g_1=\sum_i\alpha_i\chi_{e_i}\), \(g_2=\sum_i\beta_i\chi_{e_i}\). Так как носители разделены (\(g_1\) живёт в \(A_1\), \(g_2\) --- в \(A_2\), множества дизъюнктны), на каждой ступеньке \(e_i\) хотя бы одно из \(\alpha_i,\beta_i\) равно нулю: если \(\alpha_i>0\), то ступенька лежит в \(A_1\), тогда она не лежит в \(A_2\) и \(\beta_i=0\). Следовательно:

\begin{itemize}
\tightlist
\item
  \(g_1+g_2=\sum_i(\alpha_i+\beta_i)\chi_{e_i}\) --- снова \textbf{ступенчатая};
\item
  всюду на \(X\): \(g_1+g_2\le f\,\chi_A\) (там, где хоть одна ступенька ненулевая, работает соответствующее неравенство, а вторая функция там нулевая).
\end{itemize}

Отсюда цепочка: \[\int_{A_1}g_1\,d\mu+\int_{A_2}g_2\,d\mu\ \overset{(1)}{=}\ \int_A g_1\,d\mu+\int_A g_2\,d\mu\ \overset{(2)}{=}\ \int_A(g_1+g_2)\,d\mu\ \overset{(3)}{\le}\ \int_A f\,d\mu,\] где (1) --- интеграл ступенчатой не меняется при расширении множества интегрирования, если вне исходного множества функция нулевая (\(g_i=0\) вне \(A_i\subset A\)); (2) --- аддитивность интеграла для ступенчатых (непосредственно из формулы \(\sum_i(\alpha_i+\beta_i)\mu(e_i\cap A)\)); (3) --- монотонность свойством 1.

Получили «двухпараметрическое семейство неравенств» (параметры --- ступенчатые \(g_1,g_2\)). Зафиксируем \(g_1\) и перейдём к супремуму по \(g_2\): правая часть \(\int_A f\) не зависит от \(g_2\), поэтому \[\int_{A_1}g_1\,d\mu+\int_{A_2}f\,d\mu\ \le\ \int_A f\,d\mu.\] Теперь перейдём к супремуму по \(g_1\): \[\int_{A_1}f\,d\mu+\int_{A_2}f\,d\mu\ \le\ \int_A f\,d\mu.\]

\textbf{Замечание лектора (почему нельзя сразу для счётного разбиения).} При интегрировании по \(A_k\) ступенчатую нужно брать \textbf{свою} для каждого \(k\); из счётного числа кусков склеилась бы функция «с бесконечным числом ступенек» --- не ступенчатая в нашем смысле. Поэтому счётный случай сводится к конечному, а не доказывается напрямую.

\textbf{(≥), шаг 2: конечное разбиение.} \(A=A_1\sqcup\dots\sqcup A_N\) --- индукция по \(N\): \(A=A_1\sqcup(A_2\sqcup\dots\sqcup A_N)\), применяем шаг 1 и далее по индукции. Равенство для конечных разбиений доказано.

\textbf{(≥), шаг 3: счётное разбиение.} Представим \[A=A_1\sqcup\dots\sqcup A_N\sqcup B_N,\qquad B_N:=\bigsqcup_{k=N+1}^{\infty}A_k\] --- конечное разбиение из \(N+1\) множества («частичная сумма и остаток»). По шагу 2 (равенство!) \[\int_A f\,d\mu=\sum_{k=1}^{N}\int_{A_k}f\,d\mu+\int_{B_N}f\,d\mu.\] «Сделаем вопиющую вещь: последнее слагаемое выкинем» --- оно неотрицательно: \[\int_A f\,d\mu\ \ge\ \sum_{k=1}^{N}\int_{A_k}f\,d\mu\qquad\forall N.\] Предельный переход в неравенстве при \(N\to\infty\) (предел частичных сумм ряда): \[\int_A f\,d\mu\ \ge\ \sum_{k=1}^{\infty}\int_{A_k}f\,d\mu.\] Вместе с (≤) --- равенство.

\end{proof}

\textbf{Следствие 1 (монотонность по множеству).} Пусть \(A\subset B\) --- измеримые и \(0\le f\le g\) --- измеримые функции. Тогда \[\int_A f\,d\mu\le\int_B g\,d\mu.\]

\begin{proof}

По теореме \(\int_B f=\int_A f+\int_{B\setminus A}f\ge\int_A f\) (добавилось неотрицательное слагаемое \(\int_{B\setminus A}f\)), а по свойству 1 \(\int_B f\le\int_B g\).

\end{proof}

(Замечание лектора: условие неотрицательности существенно --- при отрицательных \(f,g\) на «добавке» \(B\setminus A\) неравенство может сломаться.)

\textbf{Следствие 2 (мера, порождённая функцией; плотность).} Пусть \(f\ge0\) измерима. Определим \[\nu(E):=\int_E f\,d\mu,\qquad E\in\mathcal A.\] Тогда \(\nu\) --- \textbf{мера} на \(\mathcal A\) (неотрицательна и счётно аддитивна --- в точности по теореме). В терминах названия параграфа: \(f\) --- \textbf{плотность} меры \(\nu\) относительно \(\mu\). Сам термин и его развитие (критерий плотности, единственность плотности ---) даются в последующих лекциях; здесь построена сама конструкция. \textbf{Требует сверки.} (в лекции 6 слово «плотность» как оформленный термин ещё не вводилось).

\textbf{Следствие 3 (суммируемый случай).} Пусть \(A=\bigsqcup_{k=1}^{\infty}A_k\) и \(f\) суммируема на \(A\). Тогда \[\int_A f\,d\mu=\sum_{k=1}^{\infty}\int_{A_k}f\,d\mu.\]

\begin{proof}

Теорема применена к \(f^+\ge0\) и \(f^-\ge0\): \[\int_A f^+=\sum_k\int_{A_k}f^+,\qquad \int_A f^-=\sum_k\int_{A_k}f^-,\] и оба ряда сходятся к конечным суммам. Сходящиеся ряды можно вычитать почленно и перегруппировывать: \[\int_A f=\int_A f^+-\int_A f^-=\sum_k\Bigl(\int_{A_k}f^+-\int_{A_k}f^-\Bigr)=\sum_k\int_{A_k}f. \]

\end{proof}
